\documentclass[a4paper,11pt]{article}
\usepackage[utf8]{inputenc}
\usepackage[T1]{fontenc}
\usepackage{geometry}
\geometry{margin=2.5cm}
\usepackage{booktabs}
\usepackage{enumitem}
\usepackage{xcolor}
\usepackage{listings}
\usepackage{hyperref}
\usepackage{parskip}

\definecolor{codegray}{gray}{0.9}
\lstset{
    basicstyle=\ttfamily\small,
    backgroundcolor=\color{codegray},
    breaklines=true,
    frame=single,
    numbers=none,
}

\title{\vspace{-2cm}GPU Performance \& Locking Analysis\\
       \large Vulkan Engine -- Perf Report \#03}
\author{}
\date{\today}

\begin{document}
\maketitle

\section{Overview}

This report analyses GPU-side performance bottlenecks, CPU--GPU
synchronisation stalls, and locking contention in the Vulkan engine. The
findings are drawn from a full source-code audit of all mutexes, fences,
semaphores, barriers, staging paths, compute dispatches, and threading
infrastructure.

\section{Critical Issues}

\subsection{VegetationRenderer::consolidateChunks() -- Per-Frame Stall}

\textbf{Location:} \texttt{vulkan/renderer/VegetationRenderer.cpp:316--346}

\textbf{Severity: CRITICAL}

This function is called inside \texttt{recordRead\-Barriers()}, which is
invoked from the main rendering path every frame. It creates a
\texttt{VkFence}, submits a \texttt{vkCmdCopyBuffer} command via
\texttt{vkQueueSubmit}, and then immediately calls
\texttt{vkWaitForFences(device, 1, \&fence, VK\_TRUE, UINT64\_MAX)} on the
\textbf{CPU}. The stall blocks until the GPU finishes the copy, which adds a
full pipeline drain to every frame where vegetation chunks have changed.

\textbf{Fix:} Reuse a single fence instead of creating a new one per call.
Signal it during the previous frame's geometry compute pass and check it
non-blockingly at the start of the next frame. The copy uploads should be
pipelined across frames so the CPU never blocks on the GPU.

\subsection{Octree Node Content -- Data Race in Worker Threads}

\textbf{Location:} \texttt{space/Octree.cpp:589--647}

\textbf{Severity: CRITICAL}

The \texttt{shapeChildren()} method spawns raw \texttt{std::thread} objects to
process child nodes in parallel. Each thread writes to fields of the shared
\texttt{OctreeNode} (\texttt{sdf[8]}, \texttt{bits}, \texttt{vertex},
\texttt{blockId}, \texttt{version}) without any mutual exclusion. The
\texttt{Octree::mutex} declared at line~34 of \texttt{Octree.hpp} is
\textbf{never acquired anywhere} in the codebase.

\textbf{Fix:} Either acquire \texttt{Octree::mutex} in critical sections, use
per-node atomics, or restructure the algorithm so each worker thread works on
independent nodes that do not alias in memory.

\subsection{New Thread Per Recursive Call}

\textbf{Location:} \texttt{space/Octree.cpp:589--647}

\textbf{Severity: CRITICAL}

\texttt{shapeChildren()} creates \texttt{std::thread} objects dynamically for
every non-leaf node that exceeds \texttt{minSize * 16}. Thread creation
(setup, TLS, stack allocation) and destruction overhead is incurred
repeatedly during recursive octree construction. With 8~children per node
and multiple levels of recursion, this can oversubscribe the CPU
dramatically.

The existing \texttt{ThreadPool} (\texttt{space/ThreadPool.hpp}) sits
unused---it is only consumed by the commented-out
\texttt{iterateParallelBFS} code path. All octree worker threads should be
submitted to the pool instead.

\subsection{Allocator Contention -- Centralised Bottleneck}

\textbf{Location:} \texttt{space/Allocator.tpp:34--48}

\textbf{Severity: CRITICAL}

The \texttt{Allocator<OctreeNode>} uses a single \texttt{std::shared\_mutex}
to protect its internal free list. Every \texttt{allocate()} call takes an
exclusive (\texttt{unique\_lock}) acquisition. When multiple worker threads
race to allocate new octree nodes (which happens constantly during
\texttt{shapeChildren}), all allocation is serialised on this single lock
point.

\texttt{deallocate()} also holds the exclusive lock and performs an
O(blocks) linear scan in debug builds to validate the pointer, extending the
hold time further.

\textbf{Fix:} Introduce per-thread slab caching (e.g.\ a thread-local free
list) with a fallback to the shared pool. This eliminates the centralised
bottleneck for the common case.

\section{High-Severity Issues}

\subsection{Synchronous Submit-and-Wait Pattern (runSingleTimeCommands)}

\textbf{Location:} \texttt{vulkan/VulkanApp.cpp:1351--1435}

\textbf{Severity: HIGH}

\texttt{runSingleTime\-Commands()} is the most widely used utility for
one-shot GPU work (buffer uploads, texture copies, meta-buffer updates).
Every call:
\begin{enumerate}
\item Creates a fence.
\item Allocates a command buffer from the shared \texttt{transientPool}
      (under \texttt{transientPoolMutex}).
\item Records the caller's commands.
\item Submits to the graphics queue (under \texttt{graphicsSubmitMutex}).
\item Calls \texttt{vkWaitForFences()} and blocks the CPU.
\item Destroys the fence and frees the command buffer.
\end{enumerate}

This pattern is used extensively in \texttt{IndirectRenderer::rebuild()},
\texttt{IndirectRenderer::upload\-Mesh\-VerticesAndIndices()},
\texttt{VegetationRenderer} instance uploads, \texttt{createDeviceLocalBuffer}
variants, and \texttt{TextureArray\-Manager} operations. Every invocation
blocks the calling thread until the GPU is done.

\textbf{Fix:} Convert hot-path callers to an async submission API
(\texttt{submit\-Command\-BufferAsync}) that returns a fence for lazy
synchronisation. Batch small uploads into a single submission.

\subsection{Per-Frame UBO Upload via Map/Unmap}

\textbf{Location:} \texttt{main.cpp:746--748}

\textbf{Severity: HIGH}

The main uniform buffer object (UBO) is uploaded every frame via
\texttt{vkMapMemory}~+~\texttt{memcpy}~+~\texttt{vkUnmapMemory}. This is
unnecessary on virtually all modern Vulkan drivers, which support persistent
mapping. The map/unmap pair adds driver overhead on every frame.

\textbf{Fix:} Map the UBO memory once at creation time and keep it
persistently mapped.

\subsection{Per-Frame Full-Resolution Depth Copies}

\textbf{Location:}
\texttt{vulkan/renderer/WaterRenderer.cpp:838},
\texttt{vulkan/renderer/WaterBackFaceRenderer.cpp:332}

\textbf{Severity: HIGH}

Both the water renderer and the back-face renderer issue
\texttt{vkCmdCopyImage} to copy the entire scene depth buffer into a
separate image every frame. At 1920~$\times$~1080 this copies \textasciitilde8
MiB per frame over the GPU bus, consuming memory bandwidth and adding a
barrier synchronisation point.

\textbf{Fix:} Use Vulkan input attachments (\texttt{VK\_ATTACHMENT\_LOAD\_OP\_LOAD})
to read the depth buffer during a subpass, or share depth via
\texttt{VK\_IMAGE\_LAYOUT\_DEPTH\_STENCIL\_READ\_ONLY\_OPTIMAL} instead of
copying.

\subsection{Legacy Barriers Instead of Synchronization2}

\textbf{Location:} 43 instances across the codebase
(\texttt{IndirectRenderer.cpp}, \texttt{VegetationRenderer.cpp},
\texttt{SolidRenderer.cpp}, \texttt{SkyRenderer.cpp},
\texttt{WaterRenderer.cpp}, \texttt{Solid360Renderer.cpp},
\texttt{SceneRenderer.cpp}, \texttt{VulkanApp.cpp})

\textbf{Severity: HIGH}

The project's \texttt{AGENTS.md} mandates {\itshape Synchronization2
everywhere}, yet 43 calls to the legacy \texttt{vkCmdPipelineBarrier()} API
remain. Two of these---
\texttt{IndirectRenderer.cpp:25} and line~1095---use
\texttt{VK\_PIPELINE\_STAGE\_ALL\_COMMANDS\_BIT} as the source stage mask,
which is a full pipeline stall on the GPU.

\textbf{Fix:} Convert all legacy barriers to
\texttt{VkImageMemoryBarrier2}/\texttt{vkCmdPipelineBarrier2} with
fine-grained stage masks.

\subsection{Per-Frame Descriptor Set Heap Allocations}

\textbf{Location:} \texttt{vulkan/renderer/SceneRenderer.cpp:901--902}

\textbf{Severity: HIGH}

\texttt{updateTexture\-DescriptorSet()} allocates
\texttt{VkDescriptorImageInfo} and \texttt{VkDescriptorBufferInfo} on the
heap (via \texttt{new}) for every descriptor write, then deletes them
immediately after \texttt{vkUpdateDescriptorSets}. This causes heap
allocation pressure in the per-frame update path.

\textbf{Fix:} Store descriptor info structures as member variables or use
a ring-buffer allocator for write data.

\subsection{ThreadPool Single-Mutex Bottleneck}

\textbf{Location:} \texttt{space/ThreadPool.hpp:33},
\texttt{space/ThreadPool.cpp:13}

\textbf{Severity: HIGH}

Every \texttt{enqueue()} and every worker's dequeue contend on the single
\texttt{queue\_mutex}. With N worker threads, all N+1 threads (producers +
consumers) serialise on this mutex. There is no work-stealing or per-thread
queue.

\textbf{Fix:} Replace with a lock-free MPMC queue (e.g.\ moodycamel
ConcurrentQueue, or a simple \texttt{std::atomic}-based queue adapted from
the literature). If the ThreadPool remains unused (only consumed by dead
code), remove it entirely.

\subsection{graphicsSubmitMutex Held During Blocking Wait}

\textbf{Location:} \texttt{vulkan/VulkanApp.cpp:1387--1421}

\textbf{Severity: HIGH}

\texttt{runSingleTime\-Commands()} acquires \texttt{graphics\-Submit\-Mutex}
before calling \texttt{vkQueueSubmit} and does not release it until after
\texttt{vkWaitForFences} returns. While the mutex is held, no other thread
can submit to the graphics queue, including the main render thread.

\textbf{Fix:} Release the queue mutex before waiting on the fence. The
submission has already been dispatched to the driver; the fence wait does
not need the mutex.

\subsection{Per-Frame Fence Wait (drawFrame)}

\textbf{Location:} \texttt{vulkan/VulkanApp.cpp:4518}

\textbf{Severity: HIGH}

The \texttt{drawFrame()} function waits on the in-flight fence at the start
of every frame. With only 2 frames in flight
(\texttt{MAX\_FRAMES\_IN\_FLIGHT = 2}), the GPU often finishes the previous
frame before the CPU is ready, so this wait usually returns immediately. The
low frame count limits the opportunity for CPU--GPU overlap.

\textbf{Fix:} Increase to 3 frames in flight and adjust the number of
offscreen render-target sets accordingly.

\subsection{No Vulkan Memory Allocator}

\textbf{Location:} ~30~direct\texttt{vkAllocateMemory} calls across the
codebase

\textbf{Severity: HIGH}

Every buffer and image allocation goes through \texttt{vkAllocateMemory}
directly. This causes driver overhead from many small kernel calls and
potential fragmentation. A Vulkan Memory Allocator (VMA) library would
batch allocations, sub-allocate from larger chunks, and reduce driver
overhead.

\section{Medium-Severity Issues}

\subsection{StagingRingBuffer Silent Allocation Failure}

\textbf{Location:} \texttt{vulkan/StagingRingBuffer.cpp:87--91}

\textbf{Severity: MEDIUM}

When the ring buffer is exhausted (total 4~MiB), \texttt{allocate()} returns
an empty handle without blocking. Callers must fall back to creating a
dedicated staging buffer (which incurs allocation + map + unmap overhead).
The \texttt{cv\_} condition variable is notified on \texttt{release()} but
\texttt{allocate()} never actually waits on it.

\textbf{Fix:} Increase the ring buffer size (e.g.\ to 32~MiB) and allow
callers to block on the condition variable, removing the need for dedicated
staging buffer fallbacks.

\subsection{Per-Frame Query Pool Readback}

\textbf{Location:} \texttt{main.cpp:657--681}

\textbf{Severity: MEDIUM}

\texttt{vkGetQueryPoolResults} is called every frame to read back GPU
timestamps. When results are not yet available, this can cause a partial
CPU stall. The query pool readback runs inside the render-loop hot path.

\textbf{Fix:} Use a ring of query pools (one per frame in flight) and defer
the readback by one frame, or make the readback asynchronous with a separate
staging buffer.

\subsection{async UniformBuffer mapping per-frame}

\textbf{Location:} \texttt{main.cpp:746--748}

\textbf{Severity: MEDIUM}

The UBO for the main uniform (view, projection, time, etc.) is mapped,
written, and unmapped every frame. Persistent mapping eliminates the
map/unmap calls.

\subsection{Std::sort Inside processPendingMeshes}

\textbf{Location:} \texttt{vulkan/renderer/SceneRenderer.cpp:1030--1035}

\textbf{Severity: MEDIUM}

The pending mesh queue is sorted by camera distance every frame with
\texttt{std::sort} (O(n log n)) while holding \texttt{pending\-MeshMutex}.
Under heavy tessellation load, the queue can grow large and this sort
consumes significant CPU time.

\textbf{Fix:} Use a priority queue or a small fixed-capacity buffer with
linear scan for the top-K items instead of full sort.

\subsection{Copy on Write Not Used for VkDescriptorSetLayout}

\textbf{Location:} \texttt{vulkan/VulkanApp.cpp:452--493}

The per-chunk vegetation generation compute dispatches create and allocate a
separate descriptor set, then free it after dispatch. These small allocation
and deallocation operations add up.

\textbf{Fix:} Use a ring of pre-allocated descriptor sets for the chunk
generation path.

\subsection{allocatePrimaryCommandBuffer Creates Pool Per Call}

\textbf{Location:} \texttt{vulkan/VulkanApp.cpp:2466--2490}

\textbf{Severity: MEDIUM}

The \texttt{allocatePrimary\-CommandBuffer()} function creates a new
command pool for every async command buffer allocation. This is
unnecessary overhead.

\textbf{Fix:} Use a pool ring (one per frame in flight) with reset.

\subsection{Kernel Thread Sanity}

\textbf{Location:} \texttt{space/IteratorHandler.cpp:35}

The worker loop in \texttt{iterateParallelBFS()} uses
\texttt{std::this\_thread::yield()} inside a busy loop. This code path is
currently dead (commented out in \texttt{Octree.cpp:800}).

\subsection{WrappedSignedDistanceFunction Cache is Thread-Unsafe}

\textbf{Location:} \texttt{sdf/WrappedSignedDistanceFunction.hpp:21--22}

The \texttt{cacheSDF} member and its \texttt{mtx} are declared, but the
cache is disabled by default (\texttt{cacheEnabled = false}) and the
\texttt{mtx} appears to never be acquired in implementations. If the cache
were enabled, concurrent reads/writes from worker threads during octree
shaping would cause data races.

\subsection{Potential ABBA Lock Ordering}

\textbf{Location:}
\texttt{vulkan/VulkanApp.cpp:2073} and \texttt{line~2009}

\texttt{submitCommandBufferAsync()} acquires \texttt{graphicsSubmitMutex}
then \texttt{pendingCmdMutex} (order A). \texttt{throttleIfTooManyPending()}
acquires \texttt{pendingCmdMutex} then calls
\texttt{waitForAllPendingCommandBuffers()} (order B). If the wait function
acquires \texttt{graphicsSubmitMutex}, this is a deadlock.

\textbf{Fix:} Verify \texttt{waitForAllPendingCommandBuffers} does not
acquire \texttt{graphicsSubmitMutex}, or use \texttt{std::lock()} to acquire
both mutexes atomically.

\section{Unused Mutexes (Dead Code)}

\begin{itemize}
\item \textbf{\texttt{Octree::mutex}} (\texttt{space/Octree.hpp:34})\\
      Declared but never acquired anywhere. The octree's internal data is
      accessed from multiple threads without any lock.
\item \textbf{\texttt{ThreadContext::mutex}} (\texttt{space/ThreadContext.hpp:35})\\
      Declared but never acquired. The caches it was intended to protect
      are thread-local anyway (each worker has its own \texttt{ThreadContext}).
\item \textbf{\texttt{WrappedSignedDistanceFunction::mtx}}
      (\texttt{sdf/WrappedSignedDistanceFunction.hpp:22})\\
      Declared but implementations do not acquire it. Cache is disabled by
      default anyway.
\item \textbf{\texttt{ThreadPool}} (\texttt{space/ThreadPool.hpp})\\
      The entire thread pool infrastructure is unused in production; its
      only consumer (\texttt{iterateParallelBFS}) is commented out.
\item \textbf{\texttt{ConcurrentQueue}} (\texttt{space/ConcurrentQueue.hpp})\\
      Only consumed by the dead \texttt{iterateParallelBFS} path.
\end{itemize}

\section{GPU Pipeline Summary}

The main rendering pipeline consists of the following stages, executed in
order each frame:

\begin{enumerate}
\item \textbf{Shadow map generation}~-- 3 cascade EVSM rendered via
      dedicated depth-only passes with GPU frustum culling.
\item \textbf{Depth pre-pass (Instance~1)}~-- solid geometry + real
      vegetation write depth. No color attachment.
\item \textbf{Color pass (Instance~2)}~-- sky, solid geometry
      (\texttt{LOAD\_OP\_LOAD} depth, no depth write), vegetation, and
      impostors (single-pass depth+color).
\item \textbf{Water pass}~-- copies the scene depth, renders water
      geometry and refraction, composites over the scene color.
\item \textbf{ImGui}~-- rendered into a separate swapchain pass.
\end{enumerate}

\subsection{Compute Dispatches}

\begin{itemize}
\item \textbf{Vegetation GPU cull}~-- one dispatch per frame
      (\texttt{VK\_PIPELINE\_BIND\_POINT\_COMPUTE}) with 64-wide workgroups.
\item \textbf{Indirect renderer GPU cull}~-- one dispatch per frame per
      renderer (solid, water, shadow cascades). 64-wide workgroups with
      72-byte push constants.
\item \textbf{Texture mixer}~-- dispatches for procedural texture
      generation on the compute queue. Separate command pool per queue.
\item \textbf{Vegetation instance generation}~-- per-chunk compute
      dispatches during chunk processing. Loops to handle
      \texttt{maxComputeWorkGroupCount} limits.
\end{itemize}

\section{Recommendations (Ordered by Impact)}

\begin{enumerate}
\item \textbf{Eliminate the \texttt{consolidateChunks()} per-frame stall.}
      Pipeline the vegetation chunk uploads across frames so the CPU never
      blocks on GPU copies.
\item \textbf{Fix the octree data race.} Acquire \texttt{Octree::mutex} (or
      equivalent) in worker threads or restructure to avoid shared mutable
      state.
\item \textbf{Replace \texttt{std::thread} with ThreadPool} in
      \texttt{shapeChildren()} to eliminate per-call thread creation
      overhead.
\item \textbf{Add per-thread slab caching to the Allocator} to remove the
      centralised \texttt{shared\_mutex} bottleneck.
\item \textbf{Release \texttt{graphicsSubmitMutex} before waiting on
      fences} in \texttt{runSingleTimeCommands}.
\item \textbf{Convert the 43 legacy barriers to Synchronization2} with
      fine-grained stage masks.
\item \textbf{Persistently map the main UBO} to avoid per-frame
      map/unmap driver calls.
\item \textbf{Replace per-frame depth copies} (water, back-face) with
      input attachments or shared depth reads.
\item \textbf{Integrate Vulkan Memory Allocator} to batch
      \texttt{vkAllocateMemory} calls.
\item \textbf{Remove dead code}: unused mutexes, unused ThreadPool,
      unused ConcurrentQueue.
\end{enumerate}

\end{document}
